\Ch{Overloading}{Overload}

\p When a single name is declared with two or more different declarations in the
same scope, the name is \textit{overloaded}. A declaration that declares an
overloaded name is called an \textit{overloaded declaration}. The set of
overloaded declarations that declare the same overloaded name are that name's
\textit{overload set}.

\p Only function and function template declarations can be overloaded; variable and type
declarations cannot be overloaded.

\Sec{Overloadable Declarations}{Overload.Decl}

\p This section specifies the cases in which a function declaration cannot be
overloaded. Any program that contains an invalid overload set is ill-formed.

\p An overload set is invalid if:
\begin{itemize}
  \item One or more declaration in the overload set only differ by return type.

\begin{HLSL}
int Yeet();
uint Yeet(); // ill-formed: decls differ only by return type
\end{HLSL}

  \item An overload set contains more than one member function declarations with
  the same \textit{parameter-type-list}, and one of those declarations is a
  \texttt{static} member function declaration (\ref{Classes.Static}).

  \begin{HLSL}
class Doggo {
  static void pet();
  void pet();              // ill-formed: cannot be overloaded with `static void pet()`

  static void wagTail();
  void wagTail();          // valid

  void bark(int);
  static void bark(int);   // ill-formed: cannot be overloaded with `void bark(int)`
};
\end{HLSL}

  \item An overload set contains more than one entry function declaration (\ref{Decl.Attr.Entry}).

\begin{HLSL}
[shader("vertex")]
void VS();
void VS(int);              // valid: only one entry point.

[shader("vertex")]
void Entry();

[shader("compute")]
void Entry(int);           // ill-formed: an overload set cannot have more than one entry function
\end{HLSL}

\item An overload set contains more than one function declaration which only
  differ in \texttt{const} parameter specifiers.
\begin{HLSL}
void G(int);
void G(const int);         // ill-formed: redeclaration of G(int)
void G(int) {}
void G(const int) {}       // ill-formed: redefinition of G(int)
\end{HLSL}

\item An overload set contains more than one function declaration which only
differ in parameters mismatching `out` and `inout`.

\begin{HLSL}
void H(int);
void H(in int);            // valid: redeclaration of H(int)
void H(inout int);         // valid: overloading between in and inout is allowed

void I(in int);
void I(out int);           // valid: overloading between in and out is allowed

void J(out int);
void J(inout int);         // ill-formed: Cannot overload based on out/inout mismatch
\end{HLSL}
\end{itemize}

\Sec{Declaration Matching}{Overload.Match}

\p Two function declarations with the same name and same parameter lists refer to
the same function if they are in the same scope (\ref{Basic.Scope}). A function
declaration in a containing scope is not the same as a function of the same
name and parameter list in a contained scope.

\Sec{Overload Resolution}{Overload.Res}

\p \textit{Overload resolution} is process by which a function call is mapped to
a the best overloaded function declaration. Overload resolution uses set of
functions called the \textit{candidate set}, and a list of expressions that
comprise the argument list for the call.

\p Overload resolution selects the function to call in the following
contexts:

\begin{itemize}
  \item invocation of a function named in a function call expression;
  \item invocation of a function call operator on a class object named in
  function call syntax;
  \item invocation of the operator referenced in an expression;
  \item invocation of a user-defined conversion function for copy-initialization
    of an object of class type; and
  \item invocation of a user-defined conversion function for initializing an
    object of non-class type from an expression of class type.
\end{itemize}

\p In each of these contexts a unique method is used to construct the overload
candidate set and argument expression list.

\Sub{Candidate Functions and Argument Lists}{Overload.Res.Sets}

\p At each point where an overload must be resolved an implementation will generate
a candidate set, containing all possible candidate functions, and the argument
list for the overloaded call. The subclauses under this clause
define the process by which the candidate set and argument lists are constructed
in each context where overload resolution occurs.

\p A candidate set may include both member and non-member functions to be resolved
against the same argument set. Member functions have an extra \textit{implicit object
parameter}, representing the object the function is called on. For the purposes
of overload resolution a static member function is considered to have an
implicit object parameter.

\p When appropriate the context shall construct an \textit{implied object argument}
referring to the object operated on.

\p Any implicit object parameter or implied object argument is always
the first parameter or argument.

\p The type of the implicit object parameter is: lvalue reference \texttt{cv T}; where \texttt{T}
is the class containing the declaration of the function, and \texttt{cv} is the
\textit{cv-qualification} on the function declaration.

\p Implied object arguments behave the same as any explicit argument except that no temporary
object shall be created to hold the argument for the implicit object parameter
except when the implicit object parameter is a member of a
\textit{cbuffer-declaration}.

\p If a candidate is a function template, candidate specializations are generated
through template argument deduction (\ref{Template.Deduction}), candidates are
then handled as any other candidate. A name may refer to more than one function
template, in which case candidate functions shall be generated for each function
template to form the full candidate set.

\SubSub{Function Call Syntax}{Overload.Res.Call}

\begin{innergrammar}
postfix-expression \terminal{(} \opt{expression-list} \terminal{)}
\end{innergrammar}

\p In a function call (\ref{Expr.Post.Call}), if the \textit{postfix-expression} names a
candidate set comprised of overloaded functions and/or function templates, the
overload resolves as to a named function (\ref{Overload.Res.Call.Named}). If the
\textit{postfix-expression} names an object of class type, the overload is resolved as a
call to an object of class type (\ref{Overload.Res.Call.Object}).

\SubSub{Call to a Named Function}{Overload.Res.Call.Named}

\p In a call to a named function, the \textit{postfix-expression} must be one of the
forms:

\begin{innergrammar}
\define{postfix-expression}\br
  postfix-expression \terminal{.} id-expression\br
  primary-expression
\end{innergrammar}

\p In the first grammar production, the \textit{id-expression} is the name of the function
to be called, and the \textit{postfix-expression} preceding the \texttt{.} operator. This
formation represents a \textit{qualified} function call, where the \textit{postfix-expression}
must be an object of class type, and the call is looked up as a member function
of that object. The function names found by that lookup constitute the overload
candidate set for the call. The argument list for the call is the
\textit{expression-list} in the call with the left operand of the \texttt{.} operator
prepended as the implicit object parameter (\ref{Overload.Res.Sets}).

\p In the second grammar production, the name is not qualified by an object
parameter, so it is an unqualified call. Name lookup occurs as normal for name
lookup of a function call (\ref{Basic.Lookup}). Function declarations found
constitute the candidate set. The candidate set for an unqualified call will be
comprised either entirely of non-member functions or entirely of member
functions of some class \texttt{T}. If the candidate set is entirely non-member
functions, the argument list is the \textit{expression-list}. If the candidate set is
entirely member functions of class \texttt{T}, the argument list is the
\textit{expression-list} with an implied object argument prepended. If the call
expression occurs in a context where an implicit object \texttt{this} is in scope and
of class type \texttt{T} or a type \texttt{T$^\prime$}, the \texttt{this} object will be the implicit object
argument, otherwise a contrived argument of type \texttt{T} will become the implied
object argument. If overload resolution selects a non-static member function,
and the implicit object argument is a contrived argument, the call is
ill-formed.

\SubSub{Call to a Object of Class Type}{Overload.Res.Call.Object}

\p If the \textit{primary-expression} in the function call expression is a class object of
type \texttt{T}, then the set of candidate functions will include the call operators of
\texttt{T} obtained by name lookup of the name \texttt{operator()} (\ref{Basic.Lookup}).

\p For this formulation, the argument list consists of the \textit{expression-list} from
the call syntax with the implicit object argument prefixed. The implicit object
argument in this form is the object produced from the evaluated
\textit{postfix-expression}.

\SubSub{Operators in Expressions}{Overload.Res.Oper}

\p Operators in an expression where no operands have a class or enumeration type
are assumed to be built-in operators and interpreted following \ref{Expr}.

\p If either operand is of class or enumeration type, overload resolution
determines which operator function or built-in operator is resolved for the
expression.

\p The table below maps unary and binary operator expressions to equivalent member
and non-member function call expressions. In the table \texttt{@} represents any
spellable operator, and \texttt{a} and \texttt{b} represent arguments to the operator.

\begin{center}
\begin{tabular}{|| c | c | c ||}
\hline
  Expression & Member function   & Non-member function \\
\hline
  \texttt{@a}   & \texttt{(a).operator@()}   & \texttt{operator@(a)}   \\
  \texttt{a@b}  & \texttt{(a).operator@(b)}  & \texttt{operator@(a,b)} \\
  \texttt{a[b]} & \texttt{(a).operator[](b)} &  \\
\hline
\end{tabular}
\end{center}

\p Given the table above, operand \texttt{a} is of type \texttt{T1}, and operand \texttt{b} is of type
\texttt{T2}. \texttt{T1} must be of class or enumeration type, and \texttt{T2} may be of any type
including \texttt{T1}. The candidate set for operator resolution is constructed by the
process:

\begin{itemize}
\item If \texttt{T1} is a class type, the result of qualified lookup of \texttt{operator@} is
  added to the set.
\item The result of unqualified lookup of \texttt{operator@} in the context of the
  expression excluding the results from step 1. If no operand is of class type,
  only functions where \texttt{T1} or \texttt{T2} are enumerations are added to the set.
\item The built-in operators of \texttt{operator@} from the operators defined in
  \ref{Overload.Builtin} are added to the set.
\end{itemize}

\p The left operand to built-in assignment operators, may not have user-defined
conversion applied to it, and no temporary may be introduced to hold it.

\p If the operator is the comma (\texttt{,}) operator, and there are no viable functions, the
built-in operator is used.

\SubSub{Initialization by Conversion}{Overload.Res.Convert}

\p An object of class type or non-class type \texttt{T} may be initialized with an
initializing expression of class type \texttt{S} invoking a user-defined conversion.

\p The candidate set of conversion functions will be the non-explicit conversion
functions of \texttt{S} and its base classes.

\p The argument list has one argument which is the initializer expression.

\Sub{Viable Functions}{Overload.Res.Viable}

\p Given the candidate set and argument expressions as determined by the
relevant context (\ref{Overload.Res.Sets}), a subset of viable functions can be
selected from the candidate set.

\p A function candidate $F(P_0 ... P_m)$ is not a viable function for a call
with argument list $A_0 ... A_n$ if:

\begin{itemize}
\item The function has fewer parameters than there are
  arguments in the argument list ($m < n$).
\item The function has more parameters than there are arguments to the
  argument list ($m > n$), and function parameters $ P_{n+1} ... P_m $ do not
  all have default arguments.
\item There is not an implicit conversion sequence that converts each argument
  $A_i$ to the type of the corresponding parameter $P_i$.
\end{itemize}

\Sub{Best Viable Function}{Overload.Res.Best}

\p For an overloaded call with arguments $A_0 ... A_n$, each viable function
$F(P_0 ... P_m)$, has a set of implicit conversion sequences $ICS_0(F) ... ICS_m(F)$
defining the conversion sequences for each argument $A_i$ to the
type of parameter $P_i$.

\p A viable function $F$ is defined to be a better function than another
viable function $f^\prime$ if for all arguments $ICS_i(F)$ is not a worse
conversion sequence than $ICS_i(f^\prime)$, and:

\begin{itemize}
\item for some argument $j$, $ICS_j(F)$ is a better conversion than
  $ICS_j(f^\prime)$ or,
\item in the context of an initialization by user-defined conversion, the
  conversion sequence from the return type of $F$ to the destination type is a
  better conversion sequence than the return type of $f^\prime$ to the destination
  type or,
\item $F$ is a non-template function and $f^\prime$ is a function template
  specialization, or
\item $F$ and $f^\prime$ are both function template specializations and $F$ is
  more specialized than $f^\prime$ according to function template partial ordering
  rules (\ref{Template.Func.Order}).
\end{itemize}

\p If there is one viable function that is a better function than all the other
viable functions, it is the selected function; otherwise the call is ill-formed.

\p If the resolved overload is a function with multiple declarations, and if at
least two of these declarations specify a default argument that made the
function viable, the program is ill-formed.

\begin{HLSL}
void F(int X = 1);
void F(float Y = 2.0f);

void Fn() {
  F(1);     // Okay.
  F(3.0f);  // Okay.
  F();      // Ill-formed.
}
\end{HLSL}

\Sub{Implicit Conversion Sequences}{Overload.ICS}

\p An \textit{implicit conversion sequence} is a sequence of conversions which
converts a source value to a prvalue of destination type. In overload resolution
the source value is the argument expression in a function call, and the
destination type is the type of the corresponding parameter of the function
being called.

\p When a parameter is a cxvalue an \textit{inverted implicit conversion sequence} is
required to convert the parameter type back to the argument type
for writing back to the argument expression lvalue. An inverted implicit
conversion sequence must be a well-formed implicit conversion sequence where the
source value is the implicit cxvalue of the parameter type, and the destination
type is the argument expression's lvalue type.

\p A well-formed implicit conversion sequence is either a \textit{standard conversion
sequence}, or a \textit{user-defined conversion sequence}.

\p In the following contexts, an implicit conversion sequence can only be a
standard conversion sequence:

\begin{itemize}
  \item Argument conversion for a user-defined conversion function.
  \item Copying a temporary for class copy-initialization.
  \item When passing an initializer-list as a single argument.
  \item Copy-initialization of a class by user-defined conversion.
\end{itemize}

\p An implicit conversion sequence models a copy-initialization unless it is an
inverted implicit conversion sequence when it models an assignment. Any
difference in top-level cv-qualification is handled by the copy-initialization
or assignment, and does not constitute a conversion.

\begin{note}
"Top-level" cv-qualification refers to the qualification of the value. This
means an parameter of type \texttt{T} can be initialized by a argument of type
\texttt{const T}. This does not mean that a parameter of type \texttt{inout T}
can be initialized with a argument of type \texttt{const T} because there is no
valid inverted conversion system to assign back to a value of type \texttt{const
T}.
\end{note}

\p When the source value type and the destination type are the same, the
implicit conversion sequence is an \textit{identity conversion}, which signifies
no conversion.

\p Only standard conversion sequences that do not create temporary objects are
valid for implicit object parameters or left operand to assignment operators.

\p If no sequence of conversions can be found to convert a source value to the
destination type, an implicit conversion sequence cannot be formed.

\p If several different sequences of conversions exist that convert the source
value to the destination type, the implicit conversion sequence is defined to be
the unique conversion sequence designated the \textit{ambiguous conversion sequence}.
For the purpose of ranking implicit conversion sequences, the
ambiguous conversion sequence is treated as a user-defined sequence that is
indistinguishable from any other user-defined conversion sequence. If overload
resolution selects a function using the ambiguous conversion sequence as the
best match for a call, the call is ill-formed.

\SubSub{Standard Conversion Sequences}{Overload.ICS.SCS}

\p The conversions that comprise a standard conversion sequence and the
composition of the sequence are defined in Chapter \ref{Conv}.

\p Each standard conversion is given a category and rank as defined in the table
below:

\begin{center}
  \begin{tabular}{|| c | c | c | c ||}
    \hline
Conversion & Category & Rank & Reference \\
\hline
No conversion  & Identity  &  & \\ \cline{1-2}\cline{4-4}
Lvalue-to-rvalue  & Lvalue Transformation  &  & \ref{Conv.LVal} \\
Array-to-pointer  & Lvalue Transformation  &  & \ref{Conv.Array} \\ \cline{1-2}
Qualification  & Qualification Adjustment  & Exact Match  & \ref{Conv.Qual} \\ \cline{1-4}
Vector Scalar splat (without conversion)  &  &  & \ref{Conv.VSplat} \\
Matrix Scalar splat (without conversion)  & Scalar Extension  & Extension  & \ref{Conv.MSplat} \\ \cline{1-4}
Integral promotion  &  &  & \ref{Conv.IConv} \& \ref{Conv.Rank.Int} \\
Floating point promotion  &  &  & \ref{Conv.FConv} \& \ref{Conv.Rank.Float} \\
Component-wise promotion  & Promotion  & Promotion  & \ref{Conv.CWise} \\ \cline{1-4}
Vector Scalar splat promotion  &   &  & \ref{Conv.VSplat} \\
Matrix Scalar splat promotion  & Scalar Extension Promotion  & Promotion Extension  & \ref{Conv.MSplat} \\ \cline{1-4}
Integral conversion  &  &  & \ref{Conv.IConv} \\
Floating point conversion  &  &  & \ref{Conv.FConv} \\
Floating-integral conversion  &  &  & \ref{Conv.FPInt} \\
Boolean conversion  &  &  & \ref{Conv.Bool} \\
Component-wise conversion  & Conversion  & Conversion  & \ref{Conv.CWise} \\ \cline{1-4}
Scalar splat conversion  &  &  & \ref{Conv.VSplat} \\
Vector Scalar splat conversion  &  & & \ref{Conv.VSplat} \\
Matrix Scalar splat conversion  & Scalar Extension Conversion  & Conversion Extension  & \ref{Conv.MSplat} \\ \cline{1-4}
Vector truncation (without conversion)  &  &  & \ref{Conv.VTrunc} \\
Matrix truncation (without conversion)  & Dimensionality Reduction  & Truncation  & \ref{Conv.VTrunc} \\ \cline{1-4}
Vector truncation promotion  & &  & \ref{Conv.VTrunc} \\
Matrix truncation promotion  & Dimensionality Reduction Promotion  & Promotion Truncation  & \ref{Conv.VTrunc} \\ \cline{1-4}
Vector truncation conversion  & & & \ref{Conv.VTrunc} \\
Matrix truncation conversion  & Dimensionality Reduction Conversion  & Conversion Truncation  & \ref{Conv.VTrunc} \\ \cline{1-4}
\hline
  \end{tabular}
\end{center}

\p If a scalar splat conversion occurs in a conversion sequence where all other
conversions are \textbf{Exact Match} rank, the conversion is ranked as
\textbf{Extension}. If a scalar splat occurs in a conversion
sequence with a \textbf{Promotion} conversion, the conversion is ranked as
\textbf{Promotion Extension}. If a scalar splat occurs in a conversion
sequence with a \textbf{Conversion} conversion, the conversion is ranked as
\textbf{Conversion Extension}.

\p If a vector or matrix truncation conversion occurs in a conversion sequence where all
other conversions are \textbf{Exact Match} rank, the conversion is ranked as
\textbf{Truncation}. If a vector or matrix truncation occurs in a conversion
sequence with a \textbf{Promotion} conversion, the conversion is ranked as
\textbf{Promotion Truncation}. If a vector or matrix truncation occurs in a conversion
sequence with a \textbf{Conversion} conversion, the conversion is ranked as
\textbf{Conversion Truncation}.

\p Otherwise, the rank of a conversion sequence is determined by considering the
rank of each conversion.

\p Conversion sequence ranks are ordered from better to worse as:

\begin{itemize}
  \item \textbf{Exact Match}
  \item \textbf{Extension}
  \item \textbf{Promotion}
  \item \textbf{Promotion Extension}
  \item \textbf{Conversion}
  \item \textbf{Conversion Extension}
  \item \textbf{Truncation}
  \item \textbf{Promotion Truncation}
  \item \textbf{Conversion Truncation}
\end{itemize}

\SubSub{User-defined Conversion Sequences}{Overload.ICS.User}

\p A \textit{user-defined conversion sequence} is a standard conversion sequence
followed by a user-defined conversion, followed by a second standard conversion
sequence.

\p The first standard conversion sequence converts the source type to the
type of the object parameter of the conversion function of the user-defined
conversion.

\p The second standard conversion converts the result of the user-defined
conversion to the target type for the conversion sequence.

\p If the user-defined conversion is a template specialization, the second standard
conversion sequence shall have the rank \textbf{Exact Match}.

\p If the user-defined conversion specifies a conversion from a class type to the
same class type, it shall have the rank \textbf{Exact Match} even if a conversion
function is invoked and it produces temporaries or other side-effects.

\p If the user-defined conversion converts from a class type to a base class of
that type, it shall have the rank \textbf{Conversion} even if a conversion function
is invoked and it produces temporaries or other side-effects.

\SubSub{Comparing Implicit Conversion Sequences}{Overload.ICS.Comparing}

\p A partial ordering of implicit conversion sequences exists based on defining
relationships for \textit{better conversion sequence}, and \textit{better conversion}. If an
implicit conversion sequence $ICS(f)$ is a better conversion sequence than
$ICS(f\prime)$, then the inverse is also true: $ICS(f\prime)$ is a \textit{worse conversion
sequence} than $ICS(f)$. If $ICS(f)$ is neither better nor worse than
$ICS(f\prime)$, the conversion sequences are \textit{indistinguishable conversion
sequences}.

\p A standard conversion sequence is always better than a user-defined conversion
sequence.

\p Standard conversion sequences are ordered by their ranks. Two conversion
sequences with the same rank are indistinguishable unless one of the following
rules applies:

\begin{itemize}
  \item If class \texttt{B} is derived directly or indirectly from class
  \texttt{A} and class \texttt{C} is derived directly or indirectly from class
  \texttt{B},
  \item binding of a expression of type \texttt{C} to a cxvalue of type
    \texttt{B} is better than binding an expression of type \texttt{C} to a
    cxvalue of type \texttt{A},
  \item conversion of \texttt{C} to \texttt{B} is better than conversion or
    \texttt{C} to \texttt{A},
  \item binding of a expression of type \texttt{B} to a cxvalue of type
    \texttt{A} is better than binding an expression of type \texttt{C} to a
    cxvalue of type \texttt{A},
  \item conversion of \texttt{B} to \texttt{A} is better than conversion of
    \texttt{C} to \texttt{A}.
\end{itemize}

\Sec{Overloaded Operators}{Overload.Operators}

\begin{grammar}
\define{operator-function-id}\br
  \terminal{operator} operator

\define{operator} \textnormal{one of}\br
  \terminal{+ - * / \% \^\/ \& | ~ ! < > << >> <= >= \&\& || , \(\) \[\]}
\end{grammar}

\p The operators \texttt{+}, \texttt{-}, \texttt{*}, and \texttt{\&} may only be
overloaded in their binary form.

\p The dot (\texttt{.}), \texttt{::} and \texttt{?:} operators cannot be
overloaded.

\p Operator functions may be called directly by their name in the form
\texttt{operator@} or through the resolution of an overload from an operator
expression.

\p An overloaded operator function must be either a non-static member function
or a non-member function, and must have at least one parameter whose type is a
class or enumeration.

\p An overloaded operator may not change the number of operands of an operator,
except the function call operator which may have any number of operands.

\p The meaning of the comma (\texttt{,}) operator, which is pre-defined for all
types, can be changed by operator overloading.

\p Overloaded operator functions are inherited from base classes in the same way
as other base class member functions.

\p Only an overloaded call operator may have default arguments, no other
operator functions may be defined with default arguments. All operator functions
must take exactly the number of parameters of the corresponding operator as
defined in this subclause.

\Sub{Unary Operators}{Overload.Operators.Unary}

\p An overloaded prefix unary operator must be implemented by a non-static member
function with no parameters or a non-member function with one parameter.

\p Postfix unary operators are not overloadable.

\Sub{Binary Operators}{Overload.Operators.Binary}

\p An overloaded binary operator must be implemented by a non-static member
function with one parameter or a non-member function with two parameters.

\Sub{Function Call Operator}{Overload.Operators.Call}

\p An overloaded function call operator must be a non-static member function with
an arbitrary number of parameters, and may have default arguments. Function call
operator functions are overload candidates for the function call syntax when the
\textit{postfix-expression} is an object of class type.

\Sub{Subscript Operator}{Overload.Operators.Subscript}

\p An overloaded subscript operator must be a non-static member function with one
parameter. Subscript operator functions are overload candidates for the
subscript syntax when the \textit{postfix-expression} is an object of class type.

\Sub{Built-in Operators}{Overload.Operators.BuiltIn}

\p For binary operators on arithmetic types, the result type is called the
\textit{promoted expression type} as determined by usual arithmetic conversions
(\ref{Expr.Conv}).

\p For all types \texttt{T} where \texttt{T} is an arithmetic type (\ref{Basic.Types.Arithmetic});
operators will be defined of the forms:

\begin{HLSL}
T operator+(T );
T operator-(T );
T operator~(T );
\end{HLSL}

\p For all types \texttt{L}, \texttt{R} and \texttt{LR}, where \texttt{L} and \texttt{R} are any arithmetic type
(\ref{Basic.Types.Arithmetic}), and \texttt{LR} is the promoted expression type;
operators will be defined of the forms:

\begin{HLSL}
LR operator*(L , R );
LR operator/(L , R );
LR operator+(L , R );
LR operator-(L , R );
bool operator<(L , R );
bool operator>(L , R );
bool operator<=(L , R );
bool operator>=(L , R );
bool operator==(L , R );
bool operator!=(L , R );
LR operator%(L , R );
\end{HLSL}

\p For all types \texttt{T} where \texttt{T} is an enumeration type (\ref{Decl.Enum}); operators
will be defined of the form:

\begin{HLSL}
bool operator<(T , T );
bool operator>(T , T );
bool operator<=(T , T );
bool operator>=(T , T );
bool operator==(T , T );
bool operator!=(T , T );
\end{HLSL}

\p For all types \texttt{T} where \texttt{T} is an integral type
(\ref{Basic.Types.Arithmetic}); operators will be defined of the form:

\begin{HLSL}
LR operator&(L , R );
LR operator^(L , R );
LR operator|(L , R );
L operator<<(L , R );
L operator>>(L , R );
\end{HLSL}

\p For all types \texttt{L} and \texttt{R} where \texttt{L} and \texttt{R} are arithmetic types
(\ref{Basic.Types.Arithmetic}), and \texttt{CVQ} is a \textit{cv-qualifier}; operators will be
defined of the form:

\begin{HLSL}
CVQ L & operator=(CVQ L &, R );
CVQ L & operator*=(CVQ L &, R );
CVQ L & operator/=(CVQ L &, R );
CVQ L & operator+=(CVQ L &, R );
CVQ L & operator-=(CVQ L &, R );
CVQ L & operator%=(CVQ L &, R );
\end{HLSL}

\p For all types \texttt{L} and \texttt{R} where \texttt{L} and \texttt{R} are integral types
(\ref{Basic.Types.Arithmetic}), and \texttt{CVQ} is a \textit{cv-qualifier}; operators will be
defined of the form:

\begin{HLSL}
CVQ L & operator<<=(CVQ L &, R );
CVQ L & operator>>=(CVQ L &, R );
CVQ L & operator&=(CVQ L &, R );
CVQ L & operator^=(CVQ L &, R );
CVQ L & operator|=(CVQ L &, R );
\end{HLSL}

\p For the boolean type \texttt{bool} operators will be defined of the form:

\begin{HLSL}
bool operator!(bool);
bool operator&&(bool, bool);
bool operator||(bool, bool);
\end{HLSL}

\p For all types \texttt{L}, \texttt{R} and \texttt{LR}, where \texttt{L} and \texttt{R} are any arithmetic type
(\ref{Basic.Types.Arithmetic}), and \texttt{LR} is the promoted expression type;
the ternary operator will be defined of the form:

\begin{HLSL}
LR operator?:(bool, L , R );
\end{HLSL}

\p For all types \texttt{T}, a ternary operator will be defined of the form:

\begin{HLSL}
T operator?:(bool, T , T );
\end{HLSL}
